% 01-introduction.tex — 装配源 truth: docs/paper/introduction.md（v0.3 W36 合并定稿）
% 纪律: Para1-6 六段式逐块搬运; Hero figure fig1 落 §1 尾（figure-specs.md 规格）; 数字禁改写。
% Abstract 已在 main.tex 前置区，本节从 Introduction 正文开始。label 沿用桩定义 sec:intro。
\section{Introduction}\label{sec:intro}

Training a working dog costs tens of thousands of dollars per animal, and more than half of candidates fail to graduate~\citep{grimm2026science}. Behind these numbers sits a measurement problem: organizations that breed and train dogs still rely on subjective human scoring of behavior. Automated behavior recognition would make training decisions objective and comparable across evaluators. The same need extends to livestock welfare monitoring, veterinary research, and pet behavior analysis.

Three difficulties make animal behavior recognition harder than its human counterpart. First, labeled data is scarce: behavior annotation requires expert time, and wild or working environments offer no laboratory control---and the public canine datasets that do exist are sensor time-series, synthetic renders, pose-only images, or sparsely-labeled video whose action labels are frame-indexed but carry no 3D skeletons (Section~\ref{ssec:rw-animal}). Second, skeletons dominate the practical signal---keypoint tracks are compact, background-invariant, and transfer across individuals---but skeleton-based recognition methods have been developed almost exclusively for humans. Third, and least explored, evaluation criteria are not static. Working-dog programs redefine behavior categories as training standards evolve; welfare protocols split or merge classes as regulations change. A system trained on a fixed label set must be re-annotated and retrained after every such revision.

Current literature leaves three gaps open (Section~\ref{sec:related}). Animal behavior pipelines are supervision-heavy and assume a fixed taxonomy (Section~\ref{ssec:rw-animal}). Self-supervised skeleton learning is mature but human-centric, and pretraining is studied separately from temporal segmentation (Section~\ref{ssec:rw-ssl}). Semi-supervised methods provide pseudo-label machinery whose components---seed-annotated anchors, prototype clustering, iterated confidence-filtered labeling---have, to our knowledge, never been assembled for temporal-skeleton recognition (Section~\ref{ssec:rw-semi}).

We propose PSD, a physics--semantics decoupled framework for animal behavior recognition under evolving evaluation criteria, designed for low-resource settings. PSD factorizes recognition into two layers connected by a narrow interface. A physics layer $\Phi$ learns how skeletons move without any behavior labels: it combines extreme-asymmetry contrastive pretraining adapted to quadruped topology with motion-word quantization that segments continuous streams into behavior proposals. A semantic layer $\Omega$ learns what behaviors are called from a small seed budget: rule-engine coarse labels seed class anchors, prototype clustering assigns pseudo-labels under confidence filtering, and iterated expansion operates within an annotation budget of 100--200 clips; a warm-start protocol initializes this stage from prior semantic-layer weights, making budgets as small as 20 clips usable on the synthetic-offset tier (the canine-tier outcome at a comparable budget is reported in Section~\ref{sec:experiments}). Because the layers interact only through embeddings and proposals, an evaluation-criteria transition replaces the taxonomy of $\Omega$ while $\Phi$ remains frozen.

Our contributions are fourfold:
\begin{enumerate}
\item \textbf{A physics--semantics decoupled framework} that formalizes evolving evaluation criteria as taxonomy transitions and absorbs them through semantic-layer updates alone.
\item \textbf{A transfer of image-domain anchor--cluster--pseudo-labeling to temporal-skeleton recognition that, to our knowledge, is the first of its kind}, supported by a repository-scale survey (Appendix~\ref{app:novelty}) complemented by an arXiv/CrossRef re-verification at submission time that identified point-supervised skeleton segmentation as the closest neighbor~\citep{pointsup2026}, differentiated in Section~\ref{sec:related}; we cannot exclude unpublished concurrent work.
\item \textbf{Label-free validation on public real animal-skeleton data}: self-supervised pretraining reaches 20.89\% $k$-NN top-1 against an 8.33\% random baseline, and motion-word segmentation cuts continuous streams into behavior-aligned proposals (boundary IoU \textbf{0.458 $\pm$ 0.049} versus a measured $0.32$ equal-segment-count random-cut null under a seeds pseudo-ground-truth protocol; a fixed-grid probe is statistically on par, Section~\ref{sec:ablation}).
\item \textbf{A low-resource pipeline whose end-to-end behavior we evaluate under a three-tier pre-registered protocol}, reporting its tier-dependent boundary as a finding: warm-started semantic-layer initialization makes a 20-clip budget usable (\textbf{82.0\% top-1 on the synthetic-offset tier}, the marginal budget on top of a fully labeled source domain) and the pipeline retains 90.7\% of full-budget accuracy at 10\% labels on a human-domain benchmark (a linear head on the frozen features alone retains 88.7\%), while on the canine public-real tier it is near-chance end-to-end and a generic-SSL control matches or exceeds the task-pretrained initialization at low budget on both real tiers tested (Section~\ref{sec:ablation}); the target real-K9 tier is pre-registered but unexercised. Uncertainty-based sampling is reported as an exploratory negative finding (Section~\ref{sec:ablation}).
\end{enumerate}

Across the two exercised data tiers and a human-domain verification benchmark, PSD reaches 96.6\% top-1 on the synthetic 22-class benchmark, and 82.0\% from 20 labeled clips under distribution shift (synthetic-offset tier). Run end-to-end under a corrected pseudo-label protocol, the pipeline retains 90.7\% of its full-budget linear-probe accuracy at 10\% of the labels on the human benchmark (pre-registered; clearing the 90\% line), while on the canine tier the same protocol stays near chance at a 13\% budget---a tier-dependent boundary reported as-is in Section~\ref{sec:experiments}. It attains 44.9\% on the four-class public-real subset ($1.80\times$ the 25\% random baseline; severe class imbalance disclosed in Section~\ref{sec:conclusion}), and absorbs a taxonomy transition at $\geq 3\times$ lower wall-clock retraining cost while matching full-retraining accuracy within statistical noise. Figure~\ref{fig:framework} illustrates the framework. Section~\ref{sec:related} reviews related work. Section~\ref{sec:method} details the method. Sections~\ref{sec:experiments}--\ref{sec:ablation} report experiments and ablations, and Section~\ref{sec:conclusion} concludes with limitations.

\begin{figure}[t]
  \centering
  \includegraphics[width=\linewidth]{fig1_framework_overview.pdf}
  \caption{The PSD framework. A frozen physics layer $\Phi$ (left) turns unlabeled skeleton streams into dynamics embeddings and behavior proposals via self-supervised pretraining and motion-word quantization. A revisable semantic layer $\Omega$ (right) expands rule-engine seeds into taxonomy coverage through anchor-guided clustering with iterated confidence-filtered pseudo-labeling, consolidated by self-training and warm-started small-budget initialization. The layers communicate only through embeddings and proposals, so a taxonomy transition $\mathcal{Y}\to\mathcal{Y}'$ retrains $\Omega$ alone.}\label{fig:framework}
\end{figure}
